%==============================================================================
%  Visualizacion paso a paso de ECDH y ECDSA (Alice <-> Bob).
%  Frames drop-in para presentacion_defensa(_v2).tex  -- mismo tema/colores.
%  Compilar suelto: pdflatex ecdh_ecdsa_viz.tex
%==============================================================================
\documentclass[aspectratio=169,11pt]{beamer}
\usetheme{metropolis}
\metroset{block=fill, numbering=none}

\usepackage[utf8]{inputenc}
\usepackage[T1]{fontenc}
\usepackage[spanish,es-noshorthands]{babel}
\usepackage{amsmath,amssymb}
\usepackage{tikz}
\usetikzlibrary{positioning,arrows.meta,calc,fit,backgrounds}

\definecolor{UGRdark}{HTML}{1F2933}
\definecolor{UGRaccent}{HTML}{A6192E}
\definecolor{UGRsoft}{HTML}{3E5C76}
\definecolor{UGRgreen}{HTML}{2E7D32}
\setbeamercolor{frametitle}{bg=UGRdark}
\setbeamercolor{alerted text}{fg=UGRaccent}
\setbeamercolor{block title}{bg=UGRsoft,fg=white}

% Estilos compartidos
\tikzset{
  sec/.style={draw=UGRgreen, fill=UGRgreen!10, rounded corners, align=left, inner sep=4pt},
  pub/.style={draw=UGRsoft, fill=UGRsoft!14, rounded corners, align=left, inner sep=4pt},
  sh/.style ={draw=UGRaccent, fill=UGRaccent!10, rounded corners, align=center, inner sep=4pt},
  hdr/.style={align=center, font=\footnotesize\bfseries},
  lbl/.style={UGRaccent, font=\scriptsize\bfseries},
  flow/.style={-{Stealth}, thick, UGRaccent},
}

\begin{document}

%------------------------------------------------------------------ ECDH
\begin{frame}{ECDH: intercambio de claves, paso a paso}
  \centering
  \begin{tikzpicture}[font=\scriptsize]
    % cabeceras y lineas de vida
    \node[hdr] (A) at (0,0) {Alice};
    \node[hdr] (B) at (9,0) {Bob};
    \draw[dashed, gray!60] (A) -- (0,-6.5);
    \draw[dashed, gray!60] (B) -- (9,-6.5);

    % 0. parametros del dominio (compartidos)
    \node[sh, text width=6.2cm] (T) at (4.5,-0.95)
      {\textbf{0.} Acuerdan el dominio $T=(p,a,b,G,n,h)$ \;— $G$ es público};

    % 1. par de claves
    \node[sec, text width=3.2cm, anchor=north] (ka) at (0,-1.9)
      {\textbf{1.} genera su clave\\ $d_A \in_R [1,n{-}1]$ \;\textcolor{UGRgreen}{(privada)}\\ $Q_A = d_A\,G$ \;(pública)};
    \node[sec, text width=3.2cm, anchor=north] (kb) at (9,-1.9)
      {\textbf{1.} genera su clave\\ $d_B \in_R [1,n{-}1]$ \;\textcolor{UGRgreen}{(privada)}\\ $Q_B = d_B\,G$ \;(pública)};

    % 2. intercambio por canal publico
    \draw[flow] (1.7,-4.05) -- node[above]{$Q_A$} (7.3,-4.05);
    \draw[flow] (7.3,-4.55) -- node[below]{$Q_B$} (1.7,-4.55);
    \node[lbl] at (4.5,-4.3) {2. canal público};

    % 3. cada uno calcula el secreto
    \node[pub, text width=2.7cm, anchor=north] (sa) at (0,-5.0)
      {\textbf{3.} $S = d_A\,Q_B$};
    \node[pub, text width=2.7cm, anchor=north] (sb) at (9,-5.0)
      {\textbf{3.} $S = d_B\,Q_A$};

    % 4. secreto comun
    \node[sh, text width=7.6cm, anchor=north] (res) at (4.5,-6.0)
      {\textbf{4.} $S = d_A Q_B = d_B Q_A = d_A d_B\,G$ \;(el \alert{mismo} punto)\\[2pt]
       $\Rightarrow$ clave de sesión $=\mathrm{KDF}(x_S)$};
  \end{tikzpicture}

  \vspace{0.3em}
  {\scriptsize \textcolor{UGRgreen}{verde = secreto} \quad azul = público \quad
   $\bullet$ un intruso ve $G,Q_A,Q_B$, pero hallar $d_A,d_B$ es el \alert{ECDLP}.}
\end{frame}

%------------------------------------------------------------------ ECDSA
\begin{frame}{ECDSA: firma y verificación, paso a paso}
  \centering
  \begin{tikzpicture}[font=\scriptsize]
    % FIRMANTE
    \node[hdr, anchor=north] (Sh) at (0,0.1)
      {Firmante (Alice) — privada $d$};
    \node[sec, text width=5.0cm, anchor=north] (sign) at (0,-0.4)
      {\textbf{1.} $e = H(m)$\\
       \textbf{2.} nonce $k \in_R [1,n{-}1]$ \;\textcolor{UGRgreen}{(secreto, único)}\\
       \textbf{3.} $R = kG = (x_R,y_R)$,\; $r = x_R \bmod n$\\
       \textbf{4.} $s = k^{-1}(e + d\,r) \bmod n$\\[2pt]
       $\Rightarrow$ firma $\boxed{(r,s)}$};

    % VERIFICADOR
    \node[hdr, anchor=north] (Vh) at (8.8,0.1)
      {Verificador (Bob) — pública $Q=dG$};
    \node[pub, text width=5.0cm, anchor=north] (ver) at (6.4,-0.4)
      {\textbf{1.} $e = H(m)$,\; $w = s^{-1} \bmod n$\\
       \textbf{2.} $u_1 = e\,w$,\; $u_2 = r\,w \pmod n$\\
       \textbf{3.} $R' = u_1 G + u_2 Q$\\
       \textbf{4.} aceptar si $x_{R'} \bmod n = r$};

    % transmision
    \draw[flow] (sign.east) -- node[above, font=\scriptsize]{$m,\,(r,s)$} (ver.west);

    % correccion
    \node[sh, text width=6.4cm, anchor=north] (ok) at (4.4,-3.6)
      {\textbf{Por qué cuadra:} $k = u_1 + u_2\,d \Rightarrow R' = (u_1+u_2 d)G = kG = R$};
  \end{tikzpicture}

  \vspace{0.2em}
  \begin{block}{\footnotesize El nonce $k$ es crítico}
    \footnotesize Reutilizarlo o predecirlo $\Rightarrow$ se despeja $d$ (PlayStation 3, Bitcoin).
    Solución: nonce determinista (RFC 6979).
  \end{block}
\end{frame}

\end{document}
